\chapter{TPI applied to $\mathcal{P}(\mathbb{N})$ vs $\mathbb{N}$ }


	\section{Building the TPI for our concrete case}
	
		\noindent
		Let's see if our system of relations $r_{\theta_{k}}$ accomplished the six conditions to be a TPI.
		\\\\
		
		\noindent
		This is not a circular argument. Think twice. For not making it more large I showed you first the system of relations: PHASE 01. We tried to show the potence of a TPI construction and its implications, by themself. First we showed the core idea; the core idea for groups; and the semi-xmas-tree pattern for finite cardinals and elements spectrum. AFTER we translated it to D-pairs, and in the previous chapter we defined TPI, finally, for infinite cardinals. Now we are only seeing if our system of relations from PHASE 01 is able to accomplish the six requisites of a TPI.
		\\\\
		
		\noindent
		Now you need to remember PHASE 01 very well, and check constantly this table: SYSTEM OF RELATIONS\\
		\ref{tab: system_relations_r_theta_k}
		\\\\
		
		\noindent
		FIRST REQUISITE: we need an ORIGIN set, and a partition of it. PERFECT: $LCF_{2p}$ and its partition into $\theta_{k}$ universes. We have $\aleph_{0}$ universes.
		\\\\
		
		\noindent
		SECOND REQUISITE: we need a system of relations: the partition of Abstract\_FLJA into $\aleph_{0}$ diferent $r_{\theta_{k}}$ relations fits the second one perfectly.
		\\\\
		
		\noindent
		THIRD REQUISITE: (SNEIs x SNEIs) and Families, defined in PHASE 01, fits this requisite. SNEIs is our DOMAIN set in this TPI. WE only add $F_{\infty}$ as a special subindex; AND we have said it is the first in the order.
		\\\\
		
		\noindent
		FOURTH REQUISITE: The NWSP set of each $r_{\theta_{k}}$ creates a list of nested subsets. Each one has a Family "less".
		\\\\
		
		\noindent
		FIFTH REQUISITE: we spent some time explaining WHY ALL Families had cardinal $\aleph_{1}$, the same than SNEIs.
		\\\\
		
		\noindent
		SIXTH REQUISITE: for all possible D-pair, it has a $\gamma$ level of conflict, equal to K, for example, so it belongs to the Family $F_{k}$. WE KNOW that ALL relations $r_{\theta_{M}}$, with $M > K$, give that D-pair totally disjoint PACKs, so it CAN'T belong to NWSP in those relations. For all possible D-pair, we can say easily it is going to pass from NWSP to WSP, much before finding the end of the list of relations.\\\\
		The intersection of ALL NWSP is EMPTY.
		\\\\
		
		\noindent
		\textbf{Like we could see in the table, that summarize ALL details, cited at the beggining of this section, WE HAVE A CLEAR SEMI-XMAS-TREE pattern for SNEIs (DOMAIN) vs $LCF_{2p}$(ORIGIN). WE have "shadowed" ALL (SNEIs x SNEIs).}
		\\\\
		
		\noindent
		\textbf{SNEIs has NOT a quantity of elements bigger than the quantity of elements inside $LCF_{2p}$.}
		\\\\
		
		\noindent
		This was the last subpath, we needed to find, to show why $\mathcal{P}(\mathbb{N})$ has NOT a quantity of elements, bigger than the quantity of elements inside $\mathbb{N}$.
		\\\\
		
		\noindent
		REMEMBER the general scheme of the "divide and conquer" strategy: \ref{fig:gereral_scheme}.
		\\\\
		
	\newpage
	\section{Shit-tuation vs TPI}
	
		\noindent
		Why TPI ends in the same context than the Shit-tuation? Read again carefully, unneccesary comment 6:the Shit-tuation.
		\\\\
		
		\noindent
		(SNEIs X SNEIs) will be our "S" set now: it has cardinal $\aleph_{1}$. The subindex of the cardinal is 1.
		\\\\
		
		\noindent
		1.- CHOPS are perfectly defined: Which Families of D-pairs we "transfer" from NWSP to WSP in each relation $r_{\theta_{k}}$.
		\\\\
		
		\noindent
		We have $\aleph_{0}$ relations, totally defined, so we have $\aleph_{0}$ different CHOPS. Each CHOP is the NWSP set in its relation. A subset of (SNEIs X SNEIs).
		\\\\
		
		\noindent
		2.- We have all the relations defined, so it defines instantly all NWSP sets, and each relation use a different universe $\theta_{k}$ to do "its chop". ALL CHOPS are INDEPENDENT; can be executed in PARALELL; in an INSTANT and the order of execution is not important.
		\\\\
		
		\noindent
		3.- We have an infinite lists of sub-states of (SNEIs X SNEIs). An infinite list of different subsets NWSP. THESE ARE THE PROPERTIES of the infinite list of NWSP sets:
		\\\\
		
		\noindent
		3a.- All NWSP are nested subsets of (SNEIs x SNEIs). The next one is STRICTLY contained in the previous one (in index): the difference is ONE FAMILY. Not a SUBPACK, like in DI. SUBPACKs were made of lcfs, Families are made of D-pairs. 
		\\\\
		
		\noindent
		3b.- The difference between two consecutive NWSP sets (in index) is ONE Family, that has the same cardinal as SNEIs (our DOMAIN). This was proven in PHASE 01:all Families of D-pairs have cardinal $\aleph_{1}$.
		\\\\
		
		\noindent
		3c.- The intersection of ALL NWSP sets is the empty set.
		\\\\
		
		\noindent
		THIS MEANS we have "shadowed" perfectly all (SNEIs x SNEIs) using just disjoint subsets from a partition of $LCF_{2p}$??
		\\\\
		
		\noindent
		If S ends empty, WE have transfered ALL D-pairs, from NWSP to WSP. Have NWSP, in $\aleph_{0}$ STEPS, passed from being ALL (SNEIs X SNEIs) to become the EMTPY SET???
		\\\\
		
		\noindent
		NWSP ends empty?
	
	\newpage
	\section{TWO POSSIBLE OPTIONS:}
	
	
		\noindent
		The observation about the meaning of being able to build a semi-xmas-tree pattern, IS SO STRONG AND SOLID that all mathematicians tried to solve this in the same way: we haven't "shadowed" all Families.
		\\\\
		
		
		\subsection {NWSP does not end EMPTY (S does not end empty):}
			
			\noindent
			1.- We can IGNORE the UNDENYABLE RESULT of the infinite intersection, of all sets $NWSP_{r_{\theta_{k}}}$, being EMPTY. It has no real meaning.\\\\
			2.- THE IMPORTANT POINT IS: in each STEP, NWSP always have cardinal $\aleph_{1}$. In each step we always left $\aleph_{0}$ Families of D-pairs unsolved.\\\\
			3.- No matter that we can not name a singular surviving Family inside NWSP, after all STEPS: more beyond, there is always more Families.\\\\
			4.- NWSP never ends empty !!
			\\\\
		
		\subsection{NWSP ends TOTALLY EMPTY ( S ends EMPTY):}
		
			\noindent
			1.- WE CAN NOT IGNORE THE UNDENYABLE RESULT OF THE INFINITE INTERSECTION, of all sets NWSP.\\\\
			2.- NO MATTER that each possible NWSP always have infinite cardinal, of D-pairs ( $\aleph_{1}$ ), or Families ( $\aleph_{0}$ ).\\\\
			3.- NO matter that we don't have got, a singular NWSP, being empty in the infinite list.\\\\
			4.- After all steps, you can not name a singular D-pair, surviving inside NWSP: QRE=0.
			\\\\
			
		
		\subsection{Which answer would you choose?}
		
			\noindent
			\textbf{\textit{THIS IS EXACTLY THE OPPOSITE REACTION WE GOT IN DI. NOW THEY TRY TO DEFEND THAT UNDER A SHIT-TUATION CONTEXT "S" NEVER ENDS EMPTY.}}
			\\\\
			
			\noindent
			If you say in DI S ends empty, My MEGAPACKS end empty, and I can not apply Naive CA Theorem. But you are saying NWSP ends TOTALLY EMPTY, so we have "shadowed" successfully all Families, after all STEPS of the TPI ( semi-xmas-tree pattern).
			\\\\
			
			\newpage
			\noindent
			If you try to say NWSP never ends empty, in TPI, and that the universes, do not grows until the next infinite cardinal, you are saying my MEGAPACKS in DI never ends EMPTY, so I can use the axiom of choice to define a perfect injectivity:\\
			$$ r: \mathcal{P}(\mathbb{N}) \rightarrow \mathbb{N} $$
			\\\\
			
			\noindent
			Many years ago they defended EXACTLY the answer my TPI needs. And laugh about the option I was trying to defend: S never ends empty. NOW they look at me as if I were crazy, and proceed to defend EXACTLY the answer my DI needs, CITING, word by word, my old arguments.
			\\\\
		
		\subsection {Final conclussion:}
		
		
			\noindent
			We have four options:\\\\
			1) S(MEGAPACKS) ends empty in DI: TPI is well builded.\\\\
			2) S does NOT end empty in DI: DI is well builded.\\\\
			3) S(NWSP) ends empty in TPI: TPI is well builded and is a practical equivalent of injectivity.\\\\
			4) S does not end empty in TPI: DI is well builded and is ANOTHER possible practical equivalent.
			\\\\
			
			\noindent
			No matter what you try to defend: you can not avoid to GUARANTEE one practical equivalent of injectivity EXISTS and is well builded.
			\\\\
			
			\noindent
			\textbf{\textit{Cantor's Theorem DEFEATED by a practical counterexample. We don't know which one is correct, TPI or DI, but ONE MUST BE well builded.}}
			\\\\
			\\\\
			
			
			
			
			 

			 
				
		
		
	
	
	
	
	
	
	